\begin{helper}[Replacement-pair avoidance]
Let $\beta,v,z,P_\delta$ satisfy the explicitly stated replacement-path
hypotheses: $v\in Z$, $\beta$ is missing there under both colorings, is
free for $T_\delta$, differs from $\delta$ and $\gamma$, and
$P_\delta$ is the full simple $(\gamma,\beta)$ path from $v$ to an
outside endpoint.  Neither $u^1_\delta$ nor $u^2_\delta$ occurs in the
support of $P_\delta$.  Consequently no $\gamma$-colored boundary edge
of $X_\delta$ under $\varphi_\delta$ is incident with either vertex.
Furthermore, every $\gamma$-colored edge under $\varphi_\delta$ incident
with $U_\delta$ has no incident vertex in $Z$.
\end{helper}
\begin{proof}
Put $r=\gammaIndex(\delta)$, $\gamma=\gamma(r)$,
\[
 \varphi_\delta=
 \operatorname{swap}_{\gamma,\delta}(\varphi;v^0_\delta),
 \qquad X_\delta=V(T_\delta),
 \qquad Z=V(T)\setminus
   \{u^1_\gamma,u^2_\gamma\}.
\]
We first fix all maximum and set data used below.  The supplied
\noderef{ScheideReducedConfiguration} extends the structural configuration
of \noderef{ScheideStructuralConfiguration}; in particular, it contains
the explicitly named original globally maximum triple
$(T,e_0,\varphi)$.  On the replacement side,
\noderef{ScheideReplacementMaximum}, instantiated with the present
$G,k,e_0,\varphi,T,C,D,\delta$ and replacement data, states simultaneously
that $\varphi_\delta$ is proper, that exactly $e_0$ is uncolored, and that
$(T_\delta,e_0,\varphi_\delta)$ is globally maximum.  We use this one
replacement maximum certificate throughout the proof.  Finally, the exact
vertex-set clause of \noderef{ScheideReplacementData} gives
\[
                     X_\delta=Z\cup U_\delta,
 \qquad U_\delta=\{u^1_\delta,u^2_\delta\},                 \tag{1}
\]
and says separately that both vertices of $U_\delta$ lie outside $V(T)$.
Thus $Z\cap U_\delta=\varnothing$, and the assumed $v\in Z$ also lies in
$X_\delta$.

We next prove, retaining individual edge identities, that neither path
color crosses directly between $Z$ and $U_\delta$.  Consider first an
edge $a$ with $\varphi_\delta(a)=\gamma$ incident with a vertex of
$U_\delta$.  By \noderef{ScheideReplacementPairFixed}, the same identity
satisfies $\varphi(a)=\gamma$.  If $a$ were also incident with
$w\in Z$, its other incident vertex in $U_\delta$ would be outside
$V(T)$ by \noderef{ScheideReplacementData}, whereas $w\in V(T)$.
The endpoint-pair convention in \noderef{LooplessMultigraph}, the
incidence convention in \noderef{Incident}, and the definition
\noderef{BoundaryEdges} therefore make $a$ an original boundary identity
in $\partial_G(V(T))$.  Apply \noderef{ScheideGlobalMaximumOrder} to the
original globally maximum triple inherited above.  Its last clause says
that $V(T)$ is closed.  Unfolding that clause through
\noderef{ColoredVertexSetClosed}, the boundary identity $a$ cannot have
color $\gamma$, since \noderef{ScheideDefectiveSetup} states that
$\gamma$ is missing at $v^0_\delta\in V(T)$.  This contradiction proves
\[
 \text{no $\varphi_\delta$-$\gamma$ edge joins $Z$ to $U_\delta$.} \tag{2}
\]
Notice that this proves verbatim the final assertion of the helper.

For the second color, the original missing occurrence of $\beta$ at
$v\in Z$ and the hypothesis $\beta\ne\gamma$ meet the exact assumptions
of \noderef{ScheideReplacementInternalMissing}.  Hence there is an edge
identity $b$ such that
\[
       \varphi(b)=\beta,
       \qquad G.\operatorname{ends}(b)
                    =\{u^1_\delta,u^2_\delta\}.              \tag{3}
\]
Both ends make $b$ incident with $U_\delta$, so
\noderef{ScheideReplacementPairFixed} gives
$\varphi_\delta(b)=\beta$.  The properness supplied by the already fixed
instance of \noderef{ScheideReplacementMaximum}, interpreted using
\noderef{IsProperEdgeColoring}, makes $b$ the unique
$\varphi_\delta$-$\beta$ edge incident with either of its two ends:
a distinct such edge would be adjacent to $b$ at that end by
\noderef{AdjacentEdges} and would have the same color.  Equation (3) has
both ends in $U_\delta$, so this uniqueness proves
\[
 \text{no $\varphi_\delta$-$\beta$ edge joins $Z$ to $U_\delta$.} \tag{4}
\]

Apply \noderef{ScheideSymmetricReplacementPathOrder} to the given
$\beta,v,z,P_\delta$.  It supplies exact finite path indices
$i_0,\ldots,i_4,q_v,q_f$.  At $i_0,\ldots,i_4$ the path contains, in one
orientation or the other, the consecutive nodes
\[
 f^1_\gamma,u^1_\gamma,h,u^2_\gamma,f^2_\gamma,             \tag{5}
\]
where $h$ is a $\varphi_\delta$-$\beta$ edge with the displayed two
block vertices as its exact ends.  The same result says that
\[
 i_4<q_v,\qquad q_f=q_v+1,\qquad
 P_\delta(q_v)=v^0_\delta,\qquad P_\delta(q_f)=f^0_\delta,       \tag{6}
\]
and every vertex-node after $q_v$ lies outside $X_\delta$.

Suppose for a contradiction that a vertex of $U_\delta$ occurs in the
support of $P_\delta$, and choose its first occurrence $t$.  The starting
node is $v\in Z$, so there is a $Z$-vertex occurrence before $t$; among
the finitely many such indices choose the largest one, say $q$, and call
its vertex $w\in Z$.  By the bipartite adjacency definition in
\noderef{TwoColorIncidenceGraph}, the node at $q+1$ is an edge-node,
say $a$, and the node at $q+2$ is a vertex-node, say $x$.  Both walk
adjacencies make $a$ incident with $w$ and $x$, and its
$\varphi_\delta$-color is either $\gamma$ or $\beta$.  Simplicity, which
is part of \noderef{ScheideDefectiveChainPath}, makes $w\ne x$; hence
\noderef{LooplessMultigraph} and \noderef{Incident} identify the exact
ends of $a$ as
$\{w,x\}$.

We claim that $x\notin X_\delta$.  If $x\in Z$, then its occurrence is
still before $t$ and is later than $q$, contrary to the maximal choice of
$q$.  If $x\in U_\delta$, then $a$ directly joins $Z$ to $U_\delta$,
contradicting (2) or (4) according to its color.  Equation (1) exhausts
the vertices of $X_\delta$, proving the claim.  Thus $a$ is a boundary
identity of $X_\delta$ in the exact sense of
\noderef{BoundaryEdges}.  Applying \noderef{ScheideGlobalMaximumOrder}
to the replacement globally maximum triple fixed above shows that
$X_\delta$ is closed.  Since $\beta$ is missing at $v\in X_\delta$,
\noderef{ColoredVertexSetClosed} rules out the color $\beta$ on this
boundary identity.  Therefore
\[
                         \varphi_\delta(a)=\gamma.             \tag{7}
\]

The data $\varphi_\delta(a)=\gamma$, $w\in Z$, $x\notin X_\delta$,
and $G.\operatorname{ends}(a)=\{w,x\}$ meet exactly the quantified
classification in
\noderef{ScheideSymmetricReplacementBoundaryConfiguration}.  Hence
\[
             a\in\{f^0_\delta,f^1_\gamma,f^2_\gamma\}.        \tag{8}
\]
Suppose $a$ is one of the two deleted-block identities.  Simplicity of
the path gives that edge-node a unique index, so its occurrence at
$q+1$ is one of its occurrences in (5).  It cannot be the second edge
of (5): the immediately preceding node there is one of the deleted block
vertices $u^1_\gamma,u^2_\gamma$, whereas the node at $q$ is the vertex
$w\in Z$.  If it is the first edge, the four consecutive steps in (5)
lead through the other deleted-block edge.  That last edge has, by the
endpoint clause of \noderef{ScheideStructuralConfiguration}, one block
end and one base end.  Simplicity forbids the next node from returning to
the immediately preceding block end, so the next node is its base end.
The explicit vertex positions in the same structural configuration put
every base vertex in $Z$, while the vertex-set identity in
\noderef{ScheideReplacementData} puts the two deleted block vertices
outside $X_\delta$.  Consequently this return to $Z$ occurs after $q$
and before $t$, contradicting the maximal choice of $q$.

The only remaining case in (8) is $a=f^0_\delta$.  Again path simplicity
identifies its unique occurrence $q+1$ with the index $q_f$ in (6).
Since $q+1<t$, we have $q_v< t$.  But the node at $t$ is a vertex of
$U_\delta\subseteq X_\delta$ by (1), contradicting the final clause of
(6), which places every later vertex-node outside $X_\delta$.  All three
cases in (8) are impossible.  We have proved
\[
       \operatorname{inl}(u^1_\delta),
       \operatorname{inl}(u^2_\delta)\notin\operatorname{supp}(P_\delta).
                                                                    \tag{9}
\]

It remains to derive the asserted boundary avoidance.  Instantiate
\noderef{ScheideDefectiveChainBoundaryEdges} with the same replacement
global-maximum certificate from \noderef{ScheideReplacementMaximum},
with defective color $\gamma$ as stated by
\noderef{ScheideSymmetricReplacementBoundaryConfiguration}, free color
$\beta$, starting vertex $v\in X_\delta$, its assumed missingness, and
the supplied full path.  It says that every identity in
$\partial_{\varphi_\delta,\gamma}(X_\delta)$ occurs as an edge-node in
the support of $P_\delta$; here the colored boundary is interpreted
exactly by \noderef{BoundaryColorEdges} and \noderef{BoundaryEdges}.
If such an identity $c$ were incident with $u^1_\delta$ or
$u^2_\delta$, then its color $\gamma$ and the definition
\noderef{TwoColorIncidenceGraph} would make the corresponding vertex-node
adjacent to the edge-node of $c$.  Thus both nodes lie in the same
$(\gamma,\beta)$ component.  The full-component equality in
\noderef{ScheideDefectiveChainPath} identifies this entire component with
the support of $P_\delta$, putting the pair vertex-node on the path and
contradicting (9).  This proves the boundary assertion; (2) is the final
noncrossing assertion, so all three conclusions follow.
\end{proof}
